\section{A family of paired punctures}
\label{sec:family}

We construct a family of spheres with 48 marked points.  Near one missing
parameter value these points form 24 disjoint pairs.  The purpose of this
section is to relate the actual based monodromy of the family to an action
on a free group whose weighted Magnus degree it raises.  We also identify
the quotient obtained by killing the loops around the pairs.

\subsection{The family and its torus cover}

Put
\[
 E=\C/(\Z+\Z\omega),\qquad \omega^3=1,\quad \omega\ne1,
 \qquad \iota(z)=-z.
\]
Let $q:E\to E/\langle\iota\rangle=\PP^1$ be the quotient map.
Multiplication by $\omega$ and by $4$ induce maps $\bar\omega$ and $f$
on $\PP^1$, respectively.  Thus
\[
 q\circ[4]=f\circ q,\qquad q\circ[\omega]=\bar\omega\circ q.
\]
Our parameter curve is
\begin{equation}\label{eq:parameter-base}
 B=\PP^1\setminus
 \bigl(q(E[2])\cup q(\ker(1-\omega))\bigr).
\end{equation}
For $s\in B$, let $\mathcal A_s$ be the following subset of the fiber
sphere:
\begin{equation}\label{eq:sphere-markings}
 \mathcal A_s=\{z\in\PP^1:f(z)=\bar\omega^j s
                  \text{ for some }j\in\{0,1,2\}\}.
\end{equation}

\begin{lemma}\label{lem:family}
The complement defining $B$ consists of five points.  The sets
$\mathcal A_s$ form 48 disjoint points varying as an \emph{\'{e}tale
multisection}: locally on $B$ they are the graphs of 48 distinct
holomorphic functions.  The point $q(0)$ belongs to no $\mathcal A_s$.
If $e_0\in E[2]\setminus\{0\}$ and $s_*=q(e_0)$, then as $s\to s_*$
the markings collide in 24 pairs, each at a simple ramification point
of $f$.
\end{lemma}

\begin{proof}
The group $\ker(1-\omega)$ has order three.  Its two nonzero points are
exchanged by $\iota$, and its intersection with $E[2]$ is $\{0\}$.
This gives the five excluded values.  A fixed point of $\bar\omega$
lifts to a point satisfying $\omega z=\pm z$; since $1+\omega$ is a
unit in $\Z[\omega]$, all such values are excluded by
\eqref{eq:parameter-base}.  The three targets in
\eqref{eq:sphere-markings} are therefore distinct.  They avoid
$q(E[2])$, which contains the critical values of the degree-16 map $f$.
Also $f(q(0))=q(0)$, so $q(0)$ gives a section of the complement.

The three points $e_j=\omega^j e_0$ are the distinct nonzero elements
of $E[2]$.  A point $p\in E$ with $[4]p=e_j$ is not two-torsion.
Consequently $q$ is unramified at $p$, $[4]$ is \'etale there, and $q$
is simply ramified at $e_j$.  Thus $f$ has simple ramification at
$q(p)$.  The sixteen points above each $e_j$ are paired freely by
$\iota$, giving eight collisions for each $j$ and 24 in total.
\end{proof}

Choose henceforth $e_0\in E[2]\setminus\{0\}$ and put $s_*=q(e_0)$.
Fix $s_0\in B$ sufficiently close to $s_*$, and choose
$x_0\in q^{-1}(s_0)$ close to $e_0$.  The six points
$\pm\omega^j x_0$ lie, two at a time, in disjoint $\iota$-invariant
disks $D_j$ about $e_j$, none containing zero.  Write
\[
 T=E\setminus\{\pm\omega^j x_0:0\le j\le2\},\qquad
 F=\pi_1(T,0),\qquad \Gamma=\pi_1(B,s_0).
\]
The group $\Gamma$ is free of rank four.  In the core
$E\setminus\bigcup_j\operatorname{int}D_j$, choose handle generators
$x,y$ representing the lattice basis $1,\omega$.  Choose based
meridians $a_j,a_j^-$ about the two punctures in
$D_j$, with core access paths, so that its boundary loop is
$c_j=a_ja_j^-$.  The choices and ordering can be made so that the
surface relation is $[x,y]c_0c_1c_2=1$.  Hence $F$ is free on
\begin{equation}\label{eq:torus-free-basis}
 x,y,a_0,a_1,a_2,c_0,c_1.
\end{equation}
Give the first five generators weight one and the last two weight
two.  Let $\nu_F$ be the resulting Magnus valuation on the pro-2
completion $\widehat F$.  The surface relation gives $\nu_F(c_2)\ge2$.

The multiplication map $[4]$ restricts to an unramified cover of $T$.
Its deck group and based covering subgroup are
\[
 H=(\Z/4)^2,\qquad
 K=\ker\bigl(F\longrightarrow H\bigr),
\]
where the map records the two handle coordinates modulo four.
Thus $K$ is free of rank $1+16(7-1)=97$.  Its source is
\[
 T^{(4)}=[4]^{-1}(T),\qquad K=\pi_1(T^{(4)},0),
\]
the torus with the 96 preimages of the six punctures removed.
Quotienting $T^{(4)}$ by $\iota$ gives
\[
 S=\PP^1\setminus\mathcal A_{s_0},\qquad
 \Pi=\pi_1(S,q(0)).
\]
We identify $K$ with the based fundamental group of the source torus
through $[4]$ and denote by $p:K\to\Pi$ the map induced by $q$.
The group $\Pi$ is free of rank 47.

The two fiber groups serve different purposes. The torus group $F$ has a
small basis on which monodromy can be controlled; the sphere group $\Pi$
is the fundamental group of the geometric fiber used in the surface.
The covering subgroup $K$ will connect these two descriptions. Keeping
this connection based, rather than only up to an inner automorphism,
is essential when we impose monodromy relations.

\subsection{Actual based monodromy}

Following $x_0$ over a loop $m\in\Gamma$ ends at
$(-1)^{\epsilon(m)}x_0$; this defines a character
$\epsilon:\Gamma\to\Z/2$.  Follow also the six points
$\pm\omega^j x_0$.  Their motion extends to an isotopy $H_u$ of the
torus which fixes zero and commutes with $\iota$.  Indeed, subdivide
the path into motions in disjoint small disks, prescribe a vector
field in one disk of each sign pair, extend it equivariantly to the
other, and use bump functions supported away from zero. Every $H_u$
fixes all four points of $E[2]$ pointwise: commutation with $\iota$
preserves this discrete fixed-point set, and an isotopy starting at the
identity cannot permute it. Define the
corrected endpoint and its based action by
\begin{equation}\label{eq:corrected-monodromy}
 \iota^{\epsilon(m)}H_1,
 \qquad \varphi_m\in\Aut(F).
\end{equation}
The correction fixes each of the six punctures individually; this
is what \emph{pure} means here.

\begin{proposition}\label{prop:based-monodromy}
The actions $\varphi_m$ form a homomorphism
$\Gamma\to\Aut(F)$, are pure, and commute with $\iota_*$.
They preserve $K$.  If $\beta_m\in\Aut(\Pi)$ is the actual based
monodromy of the punctured sphere family, then
\begin{equation}\label{eq:actual-monodromy-square}
 p\circ\varphi_m|_K=\beta_m\circ p.
\end{equation}
Moreover, $p:K\to\Pi$ is surjective.
\end{proposition}

\begin{proof}
Two extensions of the same point motion differ by an isotopy fixing
the initial punctures and zero.  Their endpoint actions therefore
agree.  The disk construction also extends homotopies of parameter
paths.  With loop composition chosen to agree with composition of
the endpoint actions, uncorrected monodromy is a homomorphism.
All its maps commute with $\iota$, and $\epsilon$ is a character,
so the corrected actions also compose.

The uncorrected endpoint acts trivially on the homology of the
closed torus; the correction acts there by a sign.  It therefore
preserves the handle-coordinate kernel $K$ and lifts uniquely
through $[4]$ fixing zero.  The lift commutes with $\iota$ by
uniqueness: both compositions lift the same map and fix zero.
Requiring this fixed basepoint also excludes any nontrivial deck
translation. Write $\widetilde H_u$ for the lift of $H_u$ through
$[4]$ starting at the identity; it fixes zero and gives the actual
motion of the $96$ punctures upstairs. The zero-fixed lift of the
corrected endpoint is $\iota^{\epsilon(m)}\widetilde H_1$, because
$[4]$ commutes with sign. After quotienting by sign, this endpoint
is the endpoint induced by $\widetilde H_u$, hence gives the actual
based sphere monodromy. This proves \eqref{eq:actual-monodromy-square}.

For surjectivity, represent a loop in $S$ based at $q(0)$ by a path
which avoids the four branch values of $q$ except at its endpoints.
Its open interval lifts through the ordinary two-sheeted cover
away from those values.  Both lifted endpoints tend to zero,
the unique point above $q(0)$.  Thus the lift is a based loop in
the source punctured torus, as required.
\end{proof}

\subsection{Three invariant double covers}

We next prove that $\varphi_m-\id$ raises the weighted degree on
$\widehat F$.  The only delicate first step is its action in degree
one.  Three double covers of the moving torus complement make this
action explicit.

\begin{lemma}\label{lem:degree-one}
Every $\varphi_m$, and also $\iota_*$, acts as the identity on
the weighted degree-one quotient
$V=\widehat F/\{g:\nu_F(g)\ge2\}$.
\end{lemma}

\begin{proof}
The vector space $V$ over $\F_2$ has the five basis classes
$x,y,a_0,a_1,a_2$.  Two characters of $F$ are obtained from the
handle coordinates modulo two.  Uncorrected monodromy preserves
them, and $\iota$ preserves them modulo two.

Choose an affine coordinate on the quotient sphere with
$q(0)=\infty$, using the same letter $q$ for its pullback to $E$.
It is an even meromorphic function with a double pole at zero.
Normalize its leading coefficient to one in an additive local
coordinate $v$ at zero, so that
$v(\iota z)=-v(z)$ and $q(z)=v(z)^{-2}+O(1)$.
For $j=0,1,2$, consider over the moving torus complement the cover
\begin{equation}\label{eq:character-cover}
 Y_j^2=q(z)-\bar\omega^j s.
\end{equation}
Here $s$ and $\bar\omega^j s$ are expressed in the chosen affine
coordinate.  The right side has simple zeros precisely at the
two removed $j$-punctures.  The cover is unramified elsewhere,
including at zero: with $\eta=vY_j$, its local equation is
\[
 \eta^2=v^2(q(z)-\bar\omega^j s)=1+O(v^2).
\]
Thus the points $\eta=1$ and $\eta=-1$ over zero give two global
lifts of the fixed basepoint section throughout $B$.

Let $\chi_j:F\to\F_2$ be the character of this double cover on the
fiber over $s_0$.  A cover of the entire family with a lifted
basepoint section gives an extension of $\chi_j$ to the fundamental
group of its total space.  Conjugating a fiber loop by the section
loop does not change its value in $\F_2$.  Hence uncorrected based
monodromy preserves $\chi_j$.  Sign has the lift
\[
 (z,Y_j)\longmapsto(-z,-Y_j).
\]
This lift fixes each selected sheet over zero, because both $v$
and $Y_j$ change sign and $\eta=vY_j$ does not.  It follows that
$\iota_*$ also preserves $\chi_j$, and therefore so does the
corrected action $\varphi_m$.

The character $\chi_j$ takes value one on $a_j$ and $a_j^-$ and
zero on the other puncture meridians.  In particular it vanishes
on all $c_i$ and factors through $V$.  The three characters
$\chi_j$, together with the two handle characters, form a basis
of $V^*$: evaluation on the five displayed generators has an
invertible block triangular matrix.  Their common kernel is therefore
$\{g:\nu_F(g)\ge2\}$.  Since all five characters are preserved,
this kernel is preserved and the two induced actions on $V$ are trivial.
\end{proof}

\begin{proposition}\label{prop:monodromy-raising}
For every $m\in\Gamma$, the induced operator
$\varphi_m-\id$ on the weighted completed Magnus algebra of
$\widehat F$ raises degree by at least one.
\end{proposition}

\begin{proof}
By \Cref{lem:degree-one}, the changes of the five weight-one
generators in \eqref{eq:torus-free-basis} have valuation at least
two.  It remains to check the two weight-two boundaries.
Fix one pair, writing $a=a_j$ and $a^-=a_j^-$.
Pureness gives $\varphi_m(a)=waw^{-1}$ for some $w\in F$.
Choose $d\in F$ with $\iota_*(a)=d a^-d^{-1}$.
Commutation with $\iota_*$ yields
\[
 \varphi_m(a^-)=w'a^-(w')^{-1},\qquad
 w'=\varphi_m(d)^{-1}\iota_*(w)d.
\]
By \Cref{lem:degree-one}, $w'w^{-1}$ has valuation at least two.
Replacing $w'$ by $w$ in the conjugation of $a^-$ therefore changes
it by an element of valuation at least three.  On the other hand,
conjugating $c_j=aa^-$ by $w$ changes it by valuation at least
$1+\nu_F(c_j)\ge3$.  It follows that
\[
 \nu_F\bigl(\varphi_m(c_j)c_j^{-1}\bigr)\ge3.
\]
The degree-raising substitution rule of \Cref{lem:substitution}
now applies to the entire weighted free basis.
\end{proof}

\subsection{Pinching the pairs and passing to the sphere}

We have obtained a degree-raising action on the seven-generator
torus group. We now collapse the lifted core to obtain a free group
on the individual puncture pairs. This construction will identify
both the quotient map from $K$ and its compatibility with the
\emph{actual} sphere group $\Pi$.

Write
\[
 T_{\mathrm{core}}=E\setminus\bigcup_j\operatorname{int}D_j,
 \qquad T_{\mathrm{core}}^{(4)}=[4]^{-1}(T_{\mathrm{core}}).
\]
The lifted core $T_{\mathrm{core}}^{(4)}$ is connected, because the
handle loops in $T_{\mathrm{core}}$ map onto $H$. Its complementary disks are indexed by
$(h,j)\in H\times\{0,1,2\}$. For each $j$, choose the index-zero
disk to be the one reached by lifting the chosen access path of
$a_j$ from zero; use deck translations to label the others.
All four fixed points of sign lie in $T_{\mathrm{core}}^{(4)}$, since multiplication
by four sends them to zero, outside every $D_j$.

Collapsing $T_{\mathrm{core}}^{(4)}$ kills its fundamental group, including every
disk boundary. By van Kampen, the resulting quotient of $K$ is
free on one positive meridian for each of the $48$ disks: in each
two-puncture disk, killing the boundary makes the negative meridian
the inverse of the positive one. Let $L$ be the pro-2 completion
of this free group, with generators $a_{h,j}$. Paths in the lifted
core disappear under collapse, so these generators do not depend
on the remaining choices of core access paths. Deck translations
induce the coordinate shifts $T_1,T_2$ on $L$.

We also denote the corresponding elements of the factor $H$ in
$L\rtimes H$ by $T_1,T_2$. Define
\begin{equation}\label{eq:torus-substitution}
 \begin{aligned}
 \Phi:F&\longrightarrow L\rtimes H,\\
 x&\longmapsto T_1,& y&\longmapsto T_2,&
 a_j&\longmapsto a_{0,j},& c_0,c_1&\longmapsto1.
 \end{aligned}
\end{equation}
The shifts commute and have order four, so $c_2$ also maps to one.
On the core group, $\Phi$ is exactly the $H$-valued covering
character. It therefore kills $\pi_1(T_{\mathrm{core}}^{(4)},0)\subset K$.
A positive meridian in disk $(h,j)$ can be based using the lift
of $u a_j u^{-1}$, where $u$ is a core loop with handle coordinate
$h$; its image under $\Phi$ is exactly $a_{h,j}$. Thus
$\Phi|_K:K\to L$ is the core-collapse map just constructed,
followed by pro-2 completion. Its image contains the discrete free
group on the $a_{h,j}$ and is dense in $L$.

Sign preserves the lifted core and pairs its complementary disks
freely: their centers map to nonzero two-torsion, whereas $[4]$
kills the source two-torsion. Let $b_j\in H$ be the index of the
negative of the chosen index-zero center. Sign takes a positive
puncture in disk $(h,j)$ to the negative puncture in disk
$(b_j-h,j)$. Since both core access paths and disk boundaries
have been killed, its induced action is exactly
\begin{equation}\label{eq:pinched-involution}
 \sigma(a_{h,j})=a_{b_j-h,j}^{-1}.
\end{equation}
If the index-zero center is represented by $v/4$, with
$v\bmod(\Z+\Z\omega)=e_j$, then $b_j$ is represented, up to
the deck-coordinate convention, by $-2v$ modulo
$4(\Z+\Z\omega)$. Since $e_j$ is nonzero two-torsion,
\begin{equation}\label{eq:odd-displacement}
 b_j\bmod2H\ne0.
\end{equation}

\begin{lemma}\label{lem:pinched-sphere}
Let $\Pi_{\mathrm{pair}}$ be the quotient of $\Pi$ obtained by
killing the boundary loops of its 24 two-puncture disks.
It is free of rank 24, and its pro-2 completion is naturally
\begin{equation}\label{eq:pinched-sphere-group}
 \widehat{\Pi_{\mathrm{pair}}}
 \simeq L/\overline{\langle\!\langle
              \sigma(g)g^{-1}:g\in L
                    \rangle\!\rangle}.
\end{equation}
More precisely, let $\lambda:L\twoheadrightarrow R$ be a continuous
quotient of pro-2 groups satisfying $\lambda\sigma=\lambda$.
There is a homomorphism $\theta:\Pi\to R$ with dense image and
\begin{equation}\label{eq:pinching-square}
 \theta p=\lambda\Phi|_K.
\end{equation}
For $m\in\Gamma$, the condition
\begin{equation}\label{eq:torus-equalizer-condition}
 \lambda\Phi(\varphi_m(k))=\lambda\Phi(k)
 \qquad(k\in K)
\end{equation}
implies $\theta\beta_m=\theta$ on the entire unpinched group
$\Pi$.
\end{lemma}

\begin{proof}
Each sign-paired pair of torus disks projects to one two-puncture
disk on the sphere, and either disk maps homeomorphically to it.
The image $q(T_{\mathrm{core}}^{(4)})$ is the sphere exterior, a sphere with $24$
holes. Killing its boundary loops kills its fundamental group,
so the quotient defining $\Pi_{\mathrm{pair}}$ is the quotient
induced by collapsing this exterior. Van Kampen identifies it
with the free group on one meridian per projected disk.

The map $q$ commutes with the two core collapses, with the
basepoints in their respective cores. Upstairs there is one
cyclic generator per disk, and sign identifies them in pairs by
$a_{h,j}=a_{b_j-h,j}^{-1}$. This accounts for the entire quotient:
every fixed point of sign is already in the collapsed core, and
either disk in each remaining pair maps homeomorphically to the
corresponding sphere disk. Passing to pro-2 completions proves
\eqref{eq:pinched-sphere-group}. The based collapse diagram then
gives the dense map $\theta$ and the exact identity
\eqref{eq:pinching-square}.

Finally, \Cref{prop:based-monodromy} and
\eqref{eq:torus-equalizer-condition} give
\[
 \theta\beta_m p
 =\theta p\varphi_m|_K
 =\lambda\Phi\varphi_m|_K
 =\lambda\Phi|_K
 =\theta p.
\]
The map $p$ is surjective, so $\theta\beta_m=\theta$.
This argument does not require the monodromy to preserve the
kernel of the pair-pinching map.
\end{proof}

The maps just constructed are summarized in \Cref{fig:group-maps}.
\begin{figure}[htbp]
\centering
\begin{tikzpicture}[baseline=(current bounding box.center),>=Stealth]
 \node (K) at (0,1) {$K$};
 \node (L) at (2.8,1) {$L$};
 \node (R) at (5.6,1) {$R$};
 \node (Pi) at (0,-0.4) {$\Pi$};
 \draw[->] (K) -- node[above] {$\Phi|_K$} (L);
 \draw[->] (L) -- node[above] {$\lambda$} (R);
 \draw[->] (K) -- node[left] {$p$} (Pi);
 \draw[->] (Pi) -- node[below right] {$\theta$} (R);
\end{tikzpicture}
\caption{The actual sphere group $\Pi$ maps to every pro-$2$ quotient
$R$ of $L$ on which sign acts trivially. The triangle records
$\theta p=\lambda\Phi|_K$. The map $p$ is surjective and $\Phi|_K$
has dense image. Monodromy is compared on $K$ and $\Pi$ before it is
passed to $R$.}
\label{fig:group-maps}
\end{figure}


We conclude with two properties of the finite geometric monodromy
which will control a parameter cover.  Let $\gamma_*\in\Gamma$ be
the positively oriented local circuit around $s_*$ based at the
nearby point $s_0$.

\begin{lemma}\label{lem:finite-geometric-monodromy}
The joint action of $\Gamma$ on the 48 markings and on the sign
of $x_0$ factors through a finite 2-group.  The element
$\gamma_*$ has nontrivial sign, whereas $\gamma_*^2$ fixes every
marking and acts trivially on $\Pi_{\mathrm{pair}}$.
\end{lemma}

\begin{proof}
Choose a positive representative in the source torus for each
marking.  Within the $j$th block these representatives are indexed
by $H$.  Along a parameter loop, their endpoints are obtained by
a translation in $H$ and the common sign $(-1)^{\epsilon(m)}$.
The three blocks are preserved individually.  Thus the joint action
is contained in $H^3\rtimes\Z/2$, with sign acting by inversion;
this is a finite 2-group.

The quotient $q$ is simply ramified at $e_0$, so $\gamma_*$
exchanges the two lifts of the parameter.  By the local simple
ramification in \Cref{lem:family}, each marking pair has equation
$w^2=t$ near its collision.  Its motion under $\gamma_*$ is a
half twist and under $\gamma_*^2$ a full twist.  These motions
are supported in the disjoint pair disks and fix the exterior.
A full twist acts on a disk group by conjugation by its boundary;
after that boundary is killed, it fixes the remaining cyclic
generator.  Hence the induced action on
$\Pi_{\mathrm{pair}}$ is trivial.
\end{proof}

The groups and their geometric roles are collected below. All ranks are
free-group ranks, with pro-$2$ completion understood only in the row
for $L$.
\begin{center}
\begin{tabular}{@{}llr@{}}
Group & Geometric description & Rank \\
\hline
$\Gamma$ & Five-punctured parameter sphere & $4$ \\
$F$ & Six-punctured torus & $7$ \\
$K$ & Its degree-$16$ multiplication cover & $97$ \\
$\Pi$ & Sphere with $48$ punctures & $47$ \\
$L$ & Lifted torus after killing the core and pair boundaries & $48$ \\
$\Pi_{\mathrm{pair}}$ & Sphere after killing the $24$ pair boundaries & $24$
\end{tabular}
\end{center}
The next section keeps the presentation through $L$ in order to retain
the degree estimates coming from $F$.
